\documentclass[10pt]{article}
\usepackage[margin=20mm]{geometry}
\usepackage{amsmath,amssymb,mathtools,fontspec,unicode-math,xurl,hyperref}
\IfFontExistsTF{DejaVu Serif}{\setmainfont{DejaVu Serif}}{\setmainfont{Latin Modern Roman}}
\setmathfont{latinmodern-math.otf}
\setlength{\parindent}{0pt}\setlength{\parskip}{6pt}\sloppy
\title{Holomorphic Partial Hyperbolicity: Initial Technical Review}
\author{AI-assisted technical review}
\date{}
\begin{document}
\maketitle
Abstract: The exact v3 preprint contains two Major defects in core proof dependencies: the stated isometry/contraction dichotomy fails for an explicit holomorphic product system, and the map used to prove Theorem 1.5 is not an automorphism of its stated Lie algebra. Two further Minor corrections concern an exponential lattice and a determinant degree. Explicit counterexamples, independent symbolic calculations, and repair boundaries are given below. These findings do not prove that the main existence or rigidity theorems are false. The recovery visually checked the cited preprint evidence pages.\par
\section*{Identity and review scope}
The reviewed manuscript is Disheng Xu and Jiesong Zhang, On holomorphic partially hyperbolic systems, arXiv:2401.04310v3, 24 May 2025, 29 pages, corresponding to DOI 10.4171/JEMS/1723. The bound PDF has SHA256 prefix `577a686a7a62c8d9`. The publisher records acceptance on 20 June 2025 and publication on 21 October 2025; equivalence to that later journal version is not assumed.\par
All extracted pages, Sections 1–10, both appendices and references were read. Particular attention was paid to the quantifiers in Definitions 3.1–3.2, the use of Proposition 3.3 in Lemma 8.3, the heat-semigroup and holonomy arguments, the higher-dimensional construction, and the Lie algebra calculation underlying Theorem 1.5. The formulas in Example 10.2 were additionally checked in raw PDF extraction. Independent SymPy calculations were actually executed; their script, output and exit status are retained. No author code was run.\par
\section*{Problem 1: the stated dichotomy has neither alternative in a basic product example (Major)}
\subsection*{Exact claim and admissible example}
Definition 3.2, p. 10, requires a set \(B\subset\operatorname{supp}\mu\) which has full leafwise volume in the entire center-unstable plaque \(L^{cu}_{\mathrm{loc}}(x)\) for \(\mu\)-almost every \(x\), as well as the fiberwise contraction property. Proposition 3.3, p. 11, claims that every ergodic Gibbs \(u\)-state satisfies this definition or Definition 3.1.\par
Let \(E=\mathbb C/(\mathbb Z+i\mathbb Z)\) and set \[M=E^2\times\mathbb P^1,\qquad
 A=\begin{pmatrix}2&1\\1&1\end{pmatrix},\qquad
 f(b,z)=(Ab,z/2).\] This is a holomorphic diffeomorphism of a compact complex threefold. The center fibration is the genuine product fibration with fibers \(\mathbb P^1\). In the flat metric on \(E^2\), the stable and unstable multipliers are \[\lambda_-=(3-\sqrt5)/2<1/2,\qquad
 \lambda_+=(3+\sqrt5)/2>2.\] For the spherical metric, the derivative norm of \(z\mapsto z/2\) is \[q(r)=\frac{2(1+r^2)}{4+r^2},\qquad 1/2\le q(r)\le2.\] Consequently all the strict partial-hyperbolicity and domination inequalities hold already for one iterate.\par
Let \(\mu=m_{E^2}\times\delta_0\). Haar measure is ergodic for \(A\): on the underlying real torus the eigenvalues are \(\lambda_\pm\), each repeated twice, and none is a root of unity. Thus \(\mu\) is ergodic. Its conditional measures on strong unstable plaques are Lebesgue measures, so it is a Gibbs \(u\)-state in Section 2.4. Its support is \(E^2\times\{0\}\).\par
\subsection*{Failure of both branches}
Every subset of this support has zero four-dimensional volume inside each center-unstable plaque: its intersection is contained in a two-dimensional unstable plaque. Hence no set allowed by Definition 3.2 can have the required full volume.\par
The system is not a \(\mu\)-center isometry either. Suppose the uniformly bounded fiber metrics of Definition 3.1 existed. Let \(a(b)>0\) be the norm, in the fiber metric at \((b,0)\), of the same fixed nonzero tangent vector to \(\mathbb P^1\) at \(0\). The isometry identity gives \[a(Ab)=2a(b),\qquad a(A^n b)=2^n a(b)\] on an invariant full-measure set. The uniform upper bound on the fiber distances implies a uniform upper bound on these tangent norms, contradicting the last equality. This argument permits the metrics to depend measurably on \(b\); it does not assume a common fixed spherical metric.\par
Thus Proposition 3.3 is false with its printed definition. The counterexample meets compactness, holomorphicity, fiberedness, complex dimension, domination, ergodicity and Gibbs-state hypotheses. It does not use a non-fibered quotient or an exceptional center leaf.\par
\subsection*{Impact and repair boundary}
Lemma 3.7 on p. 12 constructs a full-fiber set outside a countable exceptional set; it never proves that this set lies in \(\operatorname{supp}\mu\). The defect propagates to the dichotomy invoked in Lemma 8.3, p. 20, a key step in the proof of Theorem 1.3. This is a Major failure of a core statement and its application, even though the necessary wording correction is short.\par
Remove the requirement \(B\subset\operatorname{supp}\mu\), or replace it by \(B\subset\pi^{-1}(\operatorname{supp}\widehat\mu)\). For the explicit example, \(B=E^2\times(\mathbb P^1\setminus\{\infty\})\) then works: all pairs in a center fiber converge to \(0\), and the omitted section has zero center volume. More generally, the construction in Lemma 3.7 and the full-volume/quasiconformal argument in Section 4 use exactly this larger set, rather than the prohibited support restriction. This identifies a coherent repair to the dependency. It neither verifies the statement as printed nor constitutes a counterexample to the final rigidity theorem, which is plainly satisfied by this product example.\par
\section*{Problem 2: the construction proving Theorem 1.5 does not preserve its Lie bracket (Major)}
\subsection*{Decisive calculation}
Example 10.2, p. 25, specifies the only nonzero brackets as \[[e_1,e_2]=[e_3,e_4]=e_5\] and defines \[Fe_1=2e_1+e_3,\quad Fe_2=2e_2+e_4,\quad
 Fe_3=e_1+e_3,\quad Fe_4=e_2+e_4,\quad
 Fe_5=e_5,\quad Fe_6=e_6.\] A Lie algebra homomorphism must satisfy \([FX,FY]=F[X,Y]\). Here \[[Fe_1,Fe_2]=5e_5\ne e_5=F[e_1,e_2],
 \qquad
 [Fe_1,Fe_4]=3e_5\ne0=F[e_1,e_4].\] Thus the first assertion of the construction is false. Commuting with the almost complex structure does not repair failure to preserve the Lie bracket. A group automorphism cannot have this linear map as its derivative; the asserted lift and quotient diffeomorphism do not follow.\par
The independent calculation checked every basis pair and found four nonzero bracket residuals (for pairs \(12,14,23,34\)), while also confirming that \(FJ=JF\). This separates the false algebra claim from the true commutation calculation.\par
\subsection*{Actual impact}
The example is explicitly the proof of Theorem 1.5, which claims a compact holomorphic partially hyperbolic threefold with five-dimensional accessibility classes. The supplied construction does not establish that theorem. This is Major regardless of the location in the examples section. It does not establish nonexistence of such systems.\par
\subsection*{A verified alternative construction}
A repair requires changing the bracket or the map and checking the lattice; one cannot retain the displayed structure unchanged. For example keep \(F\) and \(Je_1=e_2\), \(Je_3=e_4\), \(Je_5=e_6\), but replace the bracket by \[[e_1,e_2]=-2e_5,\quad [e_1,e_4]=e_5,\quad
 [e_3,e_2]=e_5,\quad [e_3,e_4]=2e_5.\] Its alternating form is nondegenerate on the horizontal four-space. In complex notation its real symmetric coefficient matrix is \[Q=\begin{pmatrix}-2&1\\1&2\end{pmatrix},\qquad
 \det Q=-5,\qquad A^TQA=Q.\] These identities make \(F\) a Lie algebra automorphism. The Nijenhuis tensor vanishes on every basis pair; this was checked independently and also follows from the type \((1,1)\) bracket.\par
Take the simply connected group with this Lie algebra and the Baker–Campbell–Hausdorff lattice \[\gamma'=\mathbb Z e_1+\cdots+\mathbb Z e_4
             +\tfrac12\mathbb Z e_5+\mathbb Z e_6,\qquad
 \Gamma'=\exp(\gamma').\] The group law \(X*Y=X+Y+\tfrac12[X,Y]\) preserves \(\gamma'\). Both \(F\) and \(F^{-1}\) preserve it. Use a left quotient \(\Gamma'\backslash G\) for the left-invariant complex structure (or consistently use the right-invariant convention). It is compact and its center leaves are two-tori. The stable and unstable horizontal spaces have multipliers \(\lambda_\pm\); the center has multiplier one. Their mixed brackets generate \(\mathbb R e_5\), because the nondegenerate \(Q\) pairs the two distinct null eigenspaces nontrivially. The resulting accessibility classes have tangent space \(E^s\oplus E^u\oplus\mathbb R e_5\), of real dimension five.\par
This is a concrete replacement establishing feasibility of the claimed phenomenon; it changes the construction in the submitted proof. It is not evidence of author adoption or journal correction. On this compact quotient the classes are compact Heisenberg quotients, not copies of the simply connected group \(H_5(\mathbb R)\).\par
\section*{Problem 3: the exponential of the stated integer lattice need not be a subgroup (Minor)}
In Example 10.1, p. 24, the paper sets \(\Gamma=\exp(\operatorname{span}_{\mathbb Z}\{e_1,\ldots,e_6\})\) in the complex Heisenberg group. With the stipulated \([e_1,e_2]=e_3\), \[\log(\exp e_1\exp e_2)=e_1+e_2+\tfrac12e_3\] in exponential coordinates of the first kind. This does not belong to the indicated integer lattice. Therefore the displayed \(\Gamma\) is not a subgroup, and its quotient is not defined as written. The same issue must be addressed in Example 10.2, in addition to its bracket failure.\par
For Example 10.1, enlarge the central Gaussian-integer lattice to half Gaussian integers, keeping the horizontal lattice. The step-two BCH law then closes, and the displayed integral determinant-one horizontal map preserves the enlarged lattice in both directions. The accessible Iwasawa example is thereby recovered with the same map and Lie algebra. This isolated lattice specification does not change the rigidity or higher-dimensional construction and is classified Minor. It is not dismissed merely because it is easy to fix.\par
\section*{Problem 4: the determinant degree is missing a factor of two (Minor)}
In Proposition 9.3, p. 23, \(W=L^{\oplus 2n}\) is correctly specified, but the next display gives \(\deg(\det W)=n\deg L\). In fact \[\det W=L^{\otimes 2n},\qquad \deg(\det W)=2n\deg L.\] For example \(n=2\) gives \(4\deg L\), not \(2\deg L\). The argument uses only strict positivity, which is unchanged. This is a local numerical error with no damage to the nontriviality or non-Kähler conclusion.\par
\section*{Other proof checks and counterevidence}
The one-complex-dimensional hypotheses are essential wherever a complex derivative is called conformal. The threefold argument uses one-dimensional stable, center and unstable distributions; it must not be generalized by the unqualified sentence on p. 7 to higher-dimensional complex bundles. The separately stated higher-dimensional examples do not use that false generalization.\par
The KAK argument in Appendix A has the correct conclusion: after normalization, an unbounded sequence in \(\operatorname{PSL}(2,\mathbb C)\) has a subsequential rank-one limit, so all points except its kernel direction coalesce. This supports the repaired contraction argument.\par
For the isometric case, the heat-semigroup intertwining identity has the correct pullback/pushforward direction. Its limiting measure can be justified without assuming that arbitrary weak convergence commutes with disintegration: for each bounded continuous \(\varphi\), apply fiberwise heat convergence and dominated convergence to the fixed base measure \(\widehat\mu\). The limit is the integral of normalized fiber volumes. These are the volumes of the invariant fiber metrics, which are equivalent to the smooth leafwise volumes.\par
The actual Corollary 4.3 of Avila–Viana was checked in the author-hosted primary paper, together with the preceding holonomy-invariance setup. Its inverse version gives \(u\)-invariance when the largest fiber exponent is nonpositive. A full-measure disintegration may be extended over its holonomy-saturated set. Thus the absence of a stated local product assumption in this particular one-sided invocation is not counted as an error. The entire external proof was not re-proved.\par
Appendix B uses Gibbs \(cu\) equivalence, not just a Gibbs \(u\) measure, when applying Fubini to pairs of center points. Lusin sets and common recurrent times yield equivariance of the almost-everywhere derivative. Conjugation by a conformal map does not amplify the operator-norm error, even though a general real-linear derivative need not commute with complex rotations. Absolute continuity then gives affine linearity, and the zero normalization removes the translation. The common-time requirement comes from Lemma B.3; the two different sequence labels in the next paragraph should not be read as permission to use unrelated times.\par
The compact elliptic-fibration argument in Section 7 correctly uses constancy of the one-dimensional moduli parameter and continuity of the real-linear holonomy: the relevant lattice automorphism ambiguity is discrete modulo the finite isotropy. The higher-dimensional moduli space is not substituted here. Section 8’s open, invariant, center-saturated set and the Cesàro construction from a positive unstable-volume complement were traced through to holomorphicity almost everywhere and then everywhere.\par
For Section 9, the four proposed holomorphic sections do form a real frame. Writing \((s_1,s_2)=(a+ib,c+id)\), their real determinant is \[(a^2+b^2+c^2+d^2)^2>0.\] The independently executed symbolic expansion confirms this; no spurious dependence issue is reported. Integer hyperbolic matrices acting among the copies preserve the translation lattice. The real-analytic continuation argument distinguishes local holomorphicity from the real-analytic product trivialization. Taking an elliptic product supplies the even dimensions, without settling the separate dimension-four problem.\par
Example 10.3’s involution is free on the total threefold because of its nontrivial half-period shift. Its sixteen exceptional center fibers arise from the \(4\cdot4\) two-torsion points in the first two elliptic factors. The quotient of the base has the stated orbifold singularities, so this example is not a counterexample under the genuine fibered hypothesis.\par
\section*{Outcome, evidence and limitations}
The exact preprint has two confirmed Major findings and two Minor findings above. The Major findings damage core proof steps; neither is claimed to refute the respective main theorem. Explicit repairs and their scope are stated separately from the authors’ arguments.\par
The arXiv history confirms v3 and describes additional Section 8 arguments. The author homepage links the arXiv and publisher records. Bounded searches of these primary sources located no separate correction, supplement or released code; this is not proof that none exists. No source-version equivalence or author adoption has been verified.\par
All source text was read in the original review. The recovery reused its hash-bound substantive report and calculations and supplied actual inspection of 22 source pages, including both appendices. The historical image-access failure remains in the preserved evidence. The identified Avila–Viana dependency was checked from retained text; no external PDF page image is claimed viewed. Scientific decision is pending, human reviews are empty, and this is initial review rather than independent acceptance.\par
9 D. Xu and J. Zhang, On holomorphic partially hyperbolic systems, arXiv:2401.04310v3 (24 May 2025). Corresponding journal DOI 10.4171/JEMS/1723. A. Avila and M. Viana, Extremal Lyapunov exponents: an invariance principle and applications, author manuscript dated 15 March 2010, Sections 4.1 and Corollary 4.3, https://w3.impa.br/\textasciitilde{}avila/extremal.pdf. Only the identified dependency passages were checked. EMS Press, article metadata, https://ems.press/journals/jems/articles/14299208. D. Xu, author homepage, https://sites.google.com/view/dishengxu/homepage. Online records consulted 27 October 2026.\par
\end{document}
